\chapter{Architecture logicielle}

Ce chapitre pr\'esente l'architecture logicielle du syst\`eme Andon, en d\'ecrivant l'organisation du code source, les modules du serveur backend et de l'application frontend, ainsi que les m\'ecanismes d'int\'e-gration.

\section{Organisation globale du code}

Le code source du syst\`eme est organis\'e en deux parties principales :

\begin{itemize}
    \item \textbf{Backend (serveur Node.js) :} r\'epertoire \texttt{server/} contenant le serveur Express, le client MQTT et les routes API.
    \item \textbf{Frontend (Angular) :} r\'epertoire \texttt{client/} contenant l'application Angular.
\end{itemize}

\figplaceholder{Organisation du code source du syst\`eme}{fig:organisation_code}

\section{Le serveur backend}

\subsection{Point d'entr\'ee : server.js}

Le fichier \texttt{server.js} constitue le point d'entr\'ee principal du serveur. Il initialise les composants suivants :

\begin{enumerate}
    \item Le serveur Express avec ses middleware (JSON, CORS, authentification).
    \item Le serveur Socket.IO pour les communications temps r\'eel.
    \item La connexion au broker MQTT (Mosquitto).
    \item La connexion \`a la base de donn\'ees PostgreSQL.
    \item Le d\'emarrage du serveur HTTP sur le port 3000.
\end{enumerate}

\subsection{Structure modulaire}

Le serveur suit une architecture modulaire :

\begin{itemize}
    \item \texttt{routes/} : d\'efinition des endpoints API (REST).
    \item \texttt{middleware/} : middleware d'authentification JWT, gestion des erreurs.
    \item \texttt{models/} : mod\`eles de donn\'ees et requ\^etes SQL.
    \item \texttt{services/} : logique m\'etier (traitement MQTT, g\'en\'eration de rapports).
    \item \texttt{config/} : configuration du serveur (port, bases de donn\'ees, MQTT).
\end{itemize}

\subsection{Connexion MQTT}

Le serveur se connecte au broker MQTT au d\'emarrage et s'abonne aux topics pertinents :

\begin{lstlisting}[language=JavaScript,caption={Connexion et abonnement MQTT},label={lst:mqtt_connect}]
const mqtt = require('mqtt');
const client = mqtt.connect('mqtt://localhost:1883', {
  clientId: 'sews-server',
  clean: true,
  reconnectPeriod: 5000
});

client.on('connect', () => {
  client.subscribe('sews/andon/+/alert');
  client.subscribe('sews/capteurs/+/status');
});
\end{lstlisting}

\subsection{Diffusion Socket.IO}

Le serveur utilise Socket.IO pour diffuser les \'ev\'enements en temps r\'eel :

\begin{lstlisting}[language=JavaScript,caption={Diffusion des \'ev\'enements via Socket.IO},label={lst:socketio}]
const io = require('socket.io')(server);

io.on('connection', (socket) => {
  socket.on('subscribe', (room) => {
    socket.join(room);
  });
});

function broadcastAndonEvent(event) {
  io.to('dashboard').emit('andon-event', event);
  io.to('ligne-' + event.ligneId)
    .emit('ligne-update', event);
}
\end{lstlisting}

\subsection{API REST}

L'API REST expose les endpoints suivants :

\begin{table}[htbp]
\centering
\renewcommand{\arraystretch}{1.3}
\begin{tabular}{l l l}
\toprule
\textbf{M\'ethode} & \textbf{Endpoint} & \textbf{Description} \\
\midrule
GET & /api/lignes & Liste des lignes \\
GET & /api/lignes/:id & D\'etail d'une ligne \\
GET & /api/evenements & Liste des \'ev\'enements \\
POST & /api/evenements & Cr\'eer un \'ev\'enement \\
PUT & /api/evenements/:id & Modifier un \'ev\'enement \\
GET & /api/rapports & Statistiques et rapports \\
POST & /api/auth/login & Authentification \\
GET & /api/users & Liste des utilisateurs \\
\bottomrule
\end{tabular}
\caption{Endpoints de l'API REST}
\label{tab:api_endpoints}
\end{table}

\section{L'application frontend Angular}

\subsection{Modules fonctionnels}

L'application Angular est structur\'ee en modules fonctionnels :

\begin{itemize}
    \item \textbf{AppModule :} module racine avec le routeur et la configuration g\'en\'erale.
    \item \textbf{DashboardModule :} tableau de bord principal avec vue d'ensemble des lignes.
    \item \textbf{LignesModule :} liste et d\'etail des lignes de production.
    \item \textbf{EvenementsModule :} historique et suivi des pannes.
    \item \textbf{RapportsModule :} g\'en\'eration de rapports et statistiques.
    \item \textbf{AdminModule :} administration des utilisateurs.
    \item \textbf{AuthModule :} authentification et gestion des sessions.
\end{itemize}

\subsection{Services Angular}

Les services assurent la communication avec le backend :

\begin{lstlisting}[language=JavaScript,caption={Service de communication temps r\'eel},label={lst:socket_service}]
@Injectable({ providedIn: 'root' })
export class SocketService {
  private socket: any;

  connect() {
    this.socket = io(environment.apiUrl);
  }

  subscribeToEvent(event: string, callback: Function) {
    this.socket.on(event, callback);
  }

  joinRoom(room: string) {
    this.socket.emit('subscribe', room);
  }
}
\end{lstlisting}

\subsection{Routing}

Le routing de l'application d\'efinit les paths suivants :

\begin{table}[htbp]
\centering
\renewcommand{\arraystretch}{1.3}
\begin{tabular}{l l}
\toprule
\textbf{Path} & \textbf{Composant} \\
\midrule
/ & DashboardComponent \\
/lignes & LignesComponent \\
/lignes/:id & LigneDetailComponent \\
/evenements & EvenementsComponent \\
/rapports & RapportsComponent \\
/admin & AdminComponent \\
/login & LoginComponent \\
\bottomrule
\end{tabular}
\caption{Routes de l'application Angular}
\label{tab:routes_angular}
\end{table}

\section{Int\'e-gration des composants}

\subsection{Flux de communication}

La communication entre les composants suit le flux suivant :

\begin{enumerate}
    \item L'ESP32 publie un message MQTT.
    \item Le broker MQTT redistribue aux abonn\'es.
    \item Le serveur Node.js re\c{c}oit le message via le client MQTT.
    \item Le serveur traite et enregistre l'\'ev\'enement.
    \item Le serveur diffuse via Socket.IO.
    \item Le frontend Angular re\c{c}oit la notification.
    \item Le composant Angular met \`a jour l'affichage.
\end{enumerate}

\subsection{S\'ecurit\'e}

La s\'ecurit\'e est assur\'ee par :

\begin{itemize}
    \item \textbf{Authentification JWT :} tokens sign\'es avec expiration.
    \item \textbf{Mots de passe hash\'es :} algorithme bcrypt.
    \item \textbf{HTTPS :} chiffrement des communications.
    \item \textbf{Validation des entr\'ees :} v\'erification c\^ot\'e serveur.
    \item \textbf{Journalisation :} logs de toutes les actions critiques.
\end{itemize}

\section{Synth\`ese du chapitre}

L'architecture logicielle du syst\`eme repose sur un serveur Node.js (Express, MQTT, Socket.IO), une base de donn\'ees PostgreSQL et une application Angular. L'int\'e-gration est assur\'ee par le protocole MQTT pour les capteurs et Socket.IO pour le temps r\'eel. La structure modulaire du code facilite la maintenance et l'\'evolution du syst\`eme.
